% SOURCE_AUTHORITY: sga5_fr_workpass.tex, lines 12404--12432 (LF-normalized unit).
% SOURCE_SHA256: 791F4EFFC5E02832D5D77ED03518C8156D6F07E4C8238B03545DB93D883FBB28
\subsubsection*{3.5}

Sean $\Lambda$ un anillo conmutativo, $A$ y $B$ dos álgebras sobre $\Lambda$, y $C=A\otimes_{\Lambda}B$. Los $C$-módulos son, por tanto, los $A$-$B$-bimódulos.

Si $\Delta$ es una subcategoría triangulada plena de $D(A)$, denotaremos por $D(C)_{\Delta}$ la subcategoría triangulada plena de $D(C)$ generada por los complejos que, considerados como objetos de $D(A)$, pertenecen a $\Delta$.

\begin{statement}{Proposición 3.6}
Con las notaciones de 3.4 y 3.5, sea $\Delta$ una subcategoría triangulada plena de $D(A)$ estable bajo sumandos directos. Si $P^\bullet$ es un objeto de $D(B)_{\parf}$ (VIII 7) y $M^\bullet$ es un objeto de $D(C)_{\Delta}$, el complejo de $A$-módulos $\RHom_B(P^\bullet,M^\bullet)\in D(A)$ definido en 4.1 es un objeto de $\Delta$. Resulta, en particular, un apareamiento de categorías trianguladas
\[
\RHom:D(B)_{\parf}\times D(C)_{\Delta}\longrightarrow \Delta,
\tag{3.6.1}
\]
es decir, un bifuntor exacto en cada variable y, por consiguiente (VIII 3), una aplicación bilineal
\[
K^{\bullet}(B)\times K^{\bullet}(D(C)_{\Delta})\xrightarrow{\operatorname{hom}} K(\Delta)
\tag{3.6.2}
\]
tal que
\[
K(a)\bigl(\cl_{K^{\bullet}(B)}(P^\bullet),\cl_{K(D(C)_{\Delta})}(M^\bullet)\bigr)
=\cl_{K(\Delta)}\bigl(\RHom_B(P^\bullet,M^\bullet)\bigr).
\tag{3.6.3}
\]
Se obtiene un enunciado análogo si se sustituye $B$ por su opuesto $B^\circ$ y $\RHom_B(P^\bullet,M^\bullet)$ por $P^\bullet\otimes_B^{\mathbf L}M^\bullet$.
\end{statement}

\begin{prooftext}[Demostración]
Se demuestra primero que $\RHom_B(P^\bullet,M^\bullet)$ es un objeto de $\Delta$ en el caso particular en que $P^\bullet$ tiene una sola componente no nula, y después se razona por inducción sobre el número de componentes no nulas de $P^\bullet$.
\end{prooftext}
